\documentclass[12pt,a4paper]{article}

% Paquetes esenciales para tipografía y formato académico
\usepackage[utf8]{inputenc}
\usepackage[T1]{fontenc}
\usepackage{lmodern} % Tipografía Latin Modern - alta legibilidad
\usepackage{microtype} % Mejoras tipográficas finas
\usepackage{geometry} % Control de márgenes
\usepackage{setspace} % Control de interlineado
\usepackage{graphicx} % Para imágenes
\usepackage{float} % Control de posición de figuras
\usepackage{caption} % Mejoras en captions
\usepackage{subcaption} % Subfiguras
\usepackage{booktabs} % Tablas profesionales
\usepackage{array} % Mejoras en tablas
\usepackage{multirow} % Celdas múltiples en tablas
\usepackage{tabularx} % Tablas con ancho ajustable
\usepackage{longtable} % Tablas multipágina
\usepackage{adjustbox} % Ajuste de cajas y tablas
\usepackage{xcolor} % Colores
\usepackage{hyperref} % Enlaces y PDF interactivo
\usepackage{fancyhdr} % Encabezados y pies de página
\usepackage{enumitem} % Listas mejoradas
\usepackage{titlesec} % Títulos personalizados

% Configuración de márgenes optimizada para formato académico
\geometry{
    a4paper,
    left=2.5cm,      % Margen izquierdo estándar académico
    right=2.5cm,     % Margen derecho
    top=2.5cm,       % Margen superior
    bottom=2.5cm,    % Margen inferior
    headheight=15pt,
    headsep=20pt,
    footskip=30pt
}

% Configuración de interlineado para mejor legibilidad
\renewcommand{\baselinestretch}{1.3} % 1.3 de interlineado (estándar académico)

% Configuración de colores para alto contraste
\definecolor{darkblue}{RGB}{0,51,102}
\definecolor{lightgray}{RGB}{245,245,245}

% Configuración de hyperref para accesibilidad
\hypersetup{
    colorlinks=true,
    linkcolor=darkblue,
    filecolor=darkblue,
    urlcolor=darkblue,
    citecolor=darkblue,
    bookmarks=true,
    bookmarksopen=true,
    bookmarksnumbered=true,
    pdfstartview=FitH
}

% Configuración de fuentes y tamaños
\setlength{\parindent}{0pt} % Sin sangría
\setlength{\parskip}{8pt} % Espacio entre párrafos

% Configuración de títulos estilo académico
\titleformat{\section}{\Large\bfseries\scshape}{\thesection}{1em}{}
\titleformat{\subsection}{\large\bfseries}{\thesubsection}{1em}{}
\titleformat{\subsubsection}{\normalsize\bfseries}{\thesubsubsection}{1em}{}

% Espaciado entre secciones
\titlespacing*{\section}{0pt}{20pt}{10pt}
\titlespacing*{\subsection}{0pt}{15pt}{8pt}
\titlespacing*{\subsubsection}{0pt}{12pt}{6pt}

% Configuración de encabezado y pie de página
\pagestyle{fancy}
\fancyhf{}
\fancyhead[L]{\small\textsc{EPRO Grupo 1}}
\fancyhead[R]{\small\textsc{Universidad Escuela Colombiana de Ingeniería}}
\fancyfoot[C]{\small\thepage}
\renewcommand{\headrulewidth}{0.4pt}
\renewcommand{\footrulewidth}{0pt}

% Comandos personalizados para consistencia
\newcommand{\sectiontitle}[1]{\section{#1}}
\newcommand{\subsectiontitle}[1]{\subsection{#1}}
\newcommand{\subsubsectiontitle}[1]{\subsubsection{#1}}
\newcommand{\tablecaption}[1]{\caption{#1}}

% Configuración para tablas
\renewcommand{\arraystretch}{1.2} % Espacio en celdas de tabla
\setlength{\tabcolsep}{8pt} % Espacio entre columnas

% Configuración para figuras
\captionsetup[figure]{font=small,labelfont=bf}
\captionsetup[table]{font=small,labelfont=bf}

% Configuración de listas
\setlist[itemize]{leftmargin=*,itemsep=5pt,topsep=5pt}
\setlist[enumerate]{leftmargin=*,itemsep=5pt,topsep=5pt}

% Comando para tablas anchas
\newcommand{\adjusttable}[1]{%
  \begin{adjustbox}{width=\textwidth,center}
    #1
  \end{adjustbox}
}

% Información del documento
\title{\textbf{ESTUDIOS DE MERCADOS:}\\
MONTAJE DE UNA EMPRESA DE RESERVAS CON PLANIFICACIÓN DE VIAJE, A\\
TRAVÉS DE UNA APLICACIÓN EN COLOMBIA}
\author{EPRO Grupo 1\\
    \small Santiago Botero García\\
    \small Juana Lozano Chaves\\
    \small Juan Pablo Salamanca Mojica\\
    \small Sofia Velandia Cifuentes\\
    \small Laura Alejandra Venegas Piraban\\
    \vspace{0.5cm}
    \small Profesor: Daniel Salazar Ferro}
\date{BOGOTÁ, COLOMBIA}

\begin{document}

% Página de título
\begin{titlepage}
    \centering
    \vspace*{2cm}
    
    {\Large\textbf{UNIVERSIDAD ESCUELA COLOMBIANA DE INGENIERÍA JULIO GARAVITO}}\\[1cm]
    
    {\Large\textbf{ESTUDIOS DE MERCADOS:}}\\[0.5cm]
    {\Large MONTAJE DE UNA EMPRESA DE RESERVAS CON PLANIFICACIÓN DE VIAJE, A}\\[0.3cm]
    {\Large TRAVÉS DE UNA APLICACIÓN EN COLOMBIA}\\[2cm]
    
    {\large\textbf{EPRO GRUPO 1}}\\[1.5cm]
    
    \textbf{Estudiantes:}\\[0.5cm]
    Juan Pablo Salamanca Mojica\\
    Juana Lozano Chaves\\
    Laura Alejandra Venegas Piraban\\
    Santiago Botero García\\
    Sofia Velandia Cifuentes\\[1.5cm]
    
    \textbf{Profesor:}\\[0.5cm]
    Daniel Salazar Ferro\\[2cm]
    
    \textbf{BOGOTÁ, COLOMBIA}
    
    \vfill
\end{titlepage}

% Tabla de contenidos
\tableofcontents
\newpage

% Contenido principal
\sectiontitle{1. Oferta}

\subsectiontitle{Estructura del mercado}

El mercado de aplicaciones de reservas y planificación de viajes en Colombia es un oligopolio dominado por pocas plataformas globales como Booking.com, Airbnb, Despegar y Tripadvisor. (8-12 millones en Colombia).

\subsectiontitle{Tipo de bien}

Es un servicio digital de consumo, es decir un bien superior.

\subsectiontitle{Tipo de demanda}

Demanda elástica (Más del 60\% de los viajeros usan internet para planificar)

\subsectiontitle{Productos sustitutos}

\begin{itemize}
    \item \textbf{Agencias de viajes tradicionales (físicas):} Ofrecen asesoría personalizada, pero con menor flexibilidad y mayor costo.
    \item \textbf{Compra directa con proveedores:} Reservas en páginas de aerolíneas, hoteles o transporte sin intermediarios.
    \item \textbf{Redes sociales y buscadores:} Uso de plataformas como Instagram o Google Maps para descubrir destinos y contactar servicios.
    \item \textbf{Recomendación informal:} Grupos, foros o conocidos (ej: viajeros que comparten experiencias sin una plataforma estructurada).
\end{itemize}

\subsectiontitle{Oferta y estrategia de comercialización (6P)}

\begin{itemize}
    \item \textbf{Producto:} Plataformas integrales de reservas (alojamiento, vuelos, experiencias).
    \item \textbf{Precio:} Modelo de comisiones (10\%).
    \item \textbf{Plaza:} Canales digitales (aplicaciones y páginas web).
    \item \textbf{Promoción:} Marketing digital, reseñas y programas de fidelización.
    \item \textbf{Personas:} Enfoque en experiencia del usuario y soporte continuo.
    \item \textbf{Procesos:} Automatización de reservas, pagos digitales y recomendaciones.
\end{itemize}

\sectiontitle{2. Análisis de Competencia}

\subsectiontitle{Competidor 1: Booking.com}

\begin{itemize}
    \item \textbf{Localización geográfica:} Global, según el informe de conectividad de MinCIT (2025), es la plataforma con mayor integración de hoteles en ciudades intermedias como Armenia y Pereira.
    \item \textbf{Volumen ofrecido:} El reporte anual de Hotelagio en el 2025 indicó que superan los 28 millones de listados, con un crecimiento del 12\% en ``alojamientos no hoteleros'' en Latam.
    \item \textbf{Estrategia de comercialización:} Utiliza el programa ``Genius''. Según Booking News (2025), su algoritmo IA ajusta descuentos en tiempo real según el historial de clics del usuario.
    \item \textbf{Personas:} Viajero Pragmático. El estudio de Travel Predictions 2026 identifica a su cliente como un usuario que busca ``confirmación inmediata'' (85\% de sus usuarios corporativos).
\end{itemize}

\subsectiontitle{Competidor 2: Airbnb}

\begin{itemize}
    \item \textbf{Localización geográfica:} Global y rural. Datos de Inside Airbnb (2026) muestran que en Colombia su oferta se ha desplazado un 22\% hacia zonas no urbanas como Guatapé y el eje cafetero.
    \item \textbf{Volumen ofrecido:} Según el reporte de AirROI del 2026, Airbnb gestiona más de 7.7 millones de anuncios activos, con una tasa de ocupación promedio del 58\% en el mercado colombiano.
    \item \textbf{Estrategia de comercialización:} Se centra en categorías para diferenciar su inventario del hotelero tradicional.
    \item \textbf{Personas:} Gen Z y Nómadas. Según el Informe de Tendencias Latam 2026, el 45\% de sus huéspedes en Colombia son nómadas digitales que buscan estancias de más de 28 días.
\end{itemize}

\subsectiontitle{Competidor 3: Despegar}

\begin{itemize}
    \item \textbf{Localización geográfica:} Según Forbes Colombia (2025), es un líder latinoamericano, operando desde su sede en Argentina con oficinas físicas en Bogotá y Medellín.
    \item \textbf{Volumen ofrecido:} Despegar en su reporte de resultados de 2025, ofrecen +300 aerolíneas y 500,000 hoteles, consolidándose como la mayor OTA (Online travel agency).
    \item \textbf{Estrategia de comercialización:} Valora Analitik (2025) nos dice que captan al cliente que busca seguridad mediante la mezcla de App e interacción humana en puntos físicos.
    \item \textbf{Personas:} Planificadores Familiares. Según Boletín Turístico (2025), su segmento principal son familias que financian viajes completos (vuelo + hotel) a cuotas.
\end{itemize}

\sectiontitle{3. Demanda}

\subsectiontitle{Segmentación del mercado}

\begin{itemize}
    \item \textbf{Viajeros corporativos:} Buscan eficiencia, confirmación inmediata y programas de lealtad.
    \item \textbf{Turistas nacionales:} Buscan precios competitivos y destinos emergentes.
    \item \textbf{Turistas internacionales:} Buscan experiencias auténticas y seguridad en transacciones.
    \item \textbf{Nómadas digitales:} Buscan estancias largas y conectividad.
    \item \textbf{Familias:} Buscan paquetes completos y opciones de financiamiento.
\end{itemize}

\subsectiontitle{Comportamiento del consumidor}

\begin{itemize}
    \item Más del 60\% de los viajeros usan internet para planificar sus viajes.
    \item Crecimiento del 19,9\% en comercio electrónico en Colombia durante 2025.
    \item Mayor sensibilidad al precio debido a factores económicos.
    Importancia creciente de reseñas y recomendaciones.
    \item Preferencia por plataformas integradas que simplifiquen el proceso.
\end{itemize}

\subsectiontitle{Tendencias del mercado}

\begin{itemize}
    \item \textbf{Digitalización:} Creciente adopción de plataformas digitales para reservas.
    \item \textbf{Sostenibilidad:} Aumento del 11,5\% en demanda de destinos con baja huella de carbono.
    \item \textbf{Experiencias locales:} Mayor interés en actividades auténticas y locales.
    \item \textbf{Flexibilidad:} Demandas de políticas de cancelación flexibles.
    \item \textbf{Personalización:} Expectativa de recomendaciones basadas en preferencias individuales.
\end{itemize}

\sectiontitle{4. Análisis FODA del Mercado}

\subsectiontitle{Fortalezas del mercado}

\begin{itemize}
    \item Crecimiento sostenido del sector turístico en Colombia.
    \item Mayor conectividad aérea en aeropuertos secundarios.
    \item Aumento en el uso de plataformas digitales para reservas.
    \item Interés creciente en destinos emergentes y experiencias auténticas.
\end{itemize}

\subsectiontitle{Oportunidades del mercado}

\begin{itemize}
    \item Proyección de 21,7 millones de viajeros para cierre de 2025.
    \item Crecimiento del 6,4\% en reservas aéreas internacionales.
    \item 684 millones de transacciones en comercio electrónico durante 2025.
    \item Digitalización de la cadena turística tradicional.
    \item Mayor demanda de destinos con baja huella de carbono.
\end{itemize}

\subsectiontitle{Debilidades del mercado}

\begin{itemize}
    \item Dependencia de plataformas internacionales.
    \item Baja penetración de soluciones locales integradas.
    \item Fragmentación de proveedores turísticos pequeños.
    \item Limitaciones tecnológicas en zonas rurales.
\end{itemize}

\subsectiontitle{Amenazas del mercado}

\begin{itemize}
    \item Dominio de OTAs globales con grandes presupuestos.
    \item Volatilidad económica que afecta el poder adquisitivo.
    \item Cambios regulatorios en el sector turístico.
    \item Competencia informal a través de redes sociales.
\end{itemize}

\sectiontitle{5. Estrategias de Mercado}

\subsectiontitle{Estrategia de diferenciación}

\begin{itemize}
    \item Integración única de reservas, planificación y red social.
    \item Enfoque en destinos emergentes y proveedores locales.
    \item Tecnología de IA para personalización de experiencias.
    \item Sistema de pagos optimizado para el mercado colombiano.
\end{itemize}

\subsectiontitle{Estrategia de precios}

\begin{itemize}
    \item Modelo de comisiones competitivo (8-10\%).
    \item Precios dinámicos basados en demanda y temporada.
    \item Descuentos por volumen y lealtad.
    \item Paquetes integrados con mejores tarifas.
\end{itemize}

\subsectiontitle{Estrategia de distribución}

\begin{itemize}
    \item Aplicación móvil nativa para iOS y Android.
    \item Plataforma web responsiva.
    \item Integración con redes sociales.
    \item Alianzas estratégicas con proveedores locales.
\end{itemize}

\subsectiontitle{Estrategia de promoción}

\begin{itemize}
    \item Marketing digital enfocado en segmentos específicos.
    \item Programa de referidos y recompensas.
    \item Contenido sobre destinos emergentes.
    \item Colaboraciones con influencers de viajes.
\end{itemize}

\end{document}
